\section{Conclusions}
\label{sec:conclusions}
\emph{In short: two inference frameworks give the same raw asymmetry, $\Araw=-0.044\pm0.009\,(\text{stat})\pm0.007\,(\text{syst})$, $5.2$ standard deviations from zero, with validated intervals; what separates them is the meaning of the numbers and the tests, not their values; a machine-learning selection (Appendix~\ref{sec:ml}) confirms the result without improving it.}

We have measured the raw charge asymmetry of $B^\pm\to K^\pm K^+K^-$ decays in the 2011 LHCb Open Data record with a common selection and a common mass model, and compared how two inference frameworks treat the same question. The frequentist profile-likelihood fit and the Bayesian posterior give $\Araw=-0.0441\pm0.0085\pm0.0072$ and $-0.0443\pm0.0085\pm0.0070$ (statistical, systematic), a nonzero asymmetry at $5.2$ standard deviations that is stable under the thirteen variants of the fit and selection. The result is compatible in sign and size with the published analysis of the same data and with its later, more precise measurements, and an asymmetry concentrated in a region of the Dalitz plane, outside the $\phi(1020)$, is reproduced a priori in the region defined by the published analysis. It is a raw asymmetry: without the control channel that corrects for the production of $B^\pm$ mesons and for the detection of $K^\pm$, it is not a measurement of $\ACP$, and its significance is not that of CP violation (Sec.~\ref{sec:comparison_published}).

Each method taught something specific.
\begin{itemize}[leftmargin=*,itemsep=2pt]
\item \emph{Frequentist.} With the other eleven parameters profiled, Wilks' theorem gives a $68\%$ interval with the nominal coverage ($68.4\pm1.5\%$ in $1000$ bootstrap replicas of the data, $0.70\pm0.01$ and $0.95\pm0.01$ in pseudo-experiments from the fitted model, where the estimator is also unbiased and reaches the minimum variance bound), and the Hesse and Minos uncertainties agree, so the likelihood is parabolic in $A$ at this sample size.
\item \emph{Bayesian.} With $\sim2\times10^4$ candidates and flat priors the posterior is the likelihood, and the credible interval coincides with the confidence interval and has the same coverage in toy experiments. The Bayesian analysis adds what the frequentist one lacks: sampling diagnostics that must be checked rather than assumed (autocorrelation time, dispersed starts, label switching of the two background slopes), a Bayes factor that depends on the prior width, and a posterior predictive check whose charge statistics, being fitted, have no power.
\item \emph{Machine learning (Appendix~\ref{sec:ml}).} Weak labels from the mass spectrum are enough to train classifiers of test AUC $\simeq0.86$, but the selection does not improve the precision on $\Araw$ in a sample with little background, and it cannot remove the misidentified decays that peak under the signal. Its value is as an independent test of the stability of $\Araw$ under a different selection and, with $K$-means, as an adaptive description of the Dalitz plane that agrees with the fixed binning.
\end{itemize}

On the systematic study, the dominant contribution is the particle-identification requirement, which a finer scan shows to be consistent with a statistical fluctuation of a nested sample; the background model is second, and the signal shape does not matter. The same variants and the same combination rule give the same uncertainty in the two methods. More generally, several statistics that look informative depend on choices: the $p$-value of a binned $\chi^2$ on the binning, the Bayes factor on the prior, the coverage on whether the pseudo-experiments come from the model or from the data.

Three limitations bound the result and indicate the next steps (Sec.~\ref{sec:limitations}): no tested model reaches $\chi^2/\mathrm{ndof}\approx1$, although $\Araw$ does not depend on the signal shape; no correction for the production and detection asymmetries is applied, so $\Araw\neq\ACP$; and the local Dalitz-plane asymmetries are counts in a region, without a background subtraction by event weights. The first of these, a model of the partially reconstructed and misidentified backgrounds, would also make the machine-learning selection usable in a fit.
